\clearpage
\setcounter{aufgabe}{29}
\einheitskopf{Einheit 5 von 5 · Prozentuale Veränderung}\label{e5}

\begin{aufgabe}{Veränderung – um oder auf? Ergänze. Rechne noch nicht mit Preisen.}
\begin{teile}
\teil Ein Preis wird um $25\,\%$ gesenkt. Er sinkt auf \leerfeld[\%] des alten Preises.
\teil Ein Preis wird um $10\,\%$ gesenkt. Er sinkt auf \leerfeld[\%] des alten Preises.
\teil Ein Preis wird auf $60\,\%$ gesenkt. Er sinkt um \leerfeld[\%].
\teil Ein Preis wird um $20\,\%$ erhöht. Er steigt auf \leerfeld[\%] des alten Preises.
\teil Ein Preis wird auf $105\,\%$ erhöht. Er steigt um \leerfeld[\%].
\end{teile}
\end{aufgabe}

\begin{aufgabe}{Veränderung – Welcher Wert ist der alte? Unterstreiche in jedem Satz den alten Wert – den Wert, der $100\,\%$ ist – und schreibe ihn auf. Rechne noch nicht.}
\begin{teile}
\teil Ein Handy kostete $260\,€$, jetzt kostet es $208\,€$. \hfill alter Wert: \leerfeld
\teil Die Miete steigt von $480\,€$ auf $528\,€$. \hfill alter Wert: \leerfeld
\teil Im letzten Jahr hatte der Verein 180 Mitglieder, heute sind es 216. \hfill alter Wert: \leerfeld
\teil Der Preis ist von $3{,}20\,€$ auf $3{,}60\,€$ gestiegen. \hfill alter Wert: \leerfeld
\teil Ein Laden hat heute 45 Kunden, gestern waren es 60. \hfill alter Wert: \leerfeld
\end{teile}
\end{aufgabe}

\begin{aufgabe}{Veränderung – neuen Wert berechnen. Berechne zuerst den Prozentwert und zähle ihn dazu oder ziehe ihn ab.}
\beispiel{\text{80\,€ um 15\,\% erhöht:}\quad 1\,\% &= 0{,}80\,€ \\ 15\,\% &= 15 \cdot 0{,}80\,€ = 12{,}00\,€ \\ \text{neuer Preis} &= 80\,€ + 12\,€ = 92\,€}
\begin{teilezwei}
\tz $200\,€$ um $10\,\%$ erhöht: \leerfeld & \tz $60\,$kg um $25\,\%$ verringert: \leerfeld \\
\tz $40\,€$ um $50\,\%$ erhöht: \leerfeld & \tz $90\,$m um $10\,\%$ verringert: \leerfeld \\
\end{teilezwei}
Rechne hier mit dem Faktor.
\begin{teile}
\teil Ein Preis von $35\,€$ steigt um $20\,\%$. Rechne mit dem Faktor $1{,}2$. \leerfeld[€]
\teil Ein Preis von $65\,€$ sinkt um $20\,\%$. Rechne mit dem Faktor $0{,}8$. \leerfeld[€]
\teil Der Eintritt in ein Schwimmbad kostet $4{,}80\,€$ und wird um $25\,\%$ erhöht. Wie viel kostet der Eintritt danach? \hfill (P10 2024) \\ \leerfeld[€]
\end{teile}
\end{aufgabe}

\begin{aufgabe}{Veränderung – in Prozent berechnen. Teile die Veränderung durch den alten Wert.}
\beispiel{\text{von 50 auf 58:}\quad 58 - 50 &= 8 \\ 8 : 50 &= 0{,}16 = 16\,\%}
\begin{teilezwei}
\tz von 200 auf 220: \leerfeld[\%] & \tz von 80 auf 60: \leerfeld[\%] \\
\tz von 40 auf 52: \leerfeld[\%] & \tz von 300 auf 285: \leerfeld[\%] \\
\end{teilezwei}
\begin{teile}
\teil Ein Brötchen kostete $0{,}55\,€$ und kostet jetzt $0{,}58\,€$. Um wie viel Prozent ist der Preis gestiegen? Runde auf eine Stelle nach dem Komma. \hfill (P10 2026) \\ \leerfeld[\%]
\end{teile}
\end{aufgabe}

\begin{aufgabe}{Veränderung – alten Wert zurückrechnen. Überlege zuerst, wie viel Prozent der neue Wert ist.}
\beispiel{\text{um 20\,\% gesenkt, jetzt 48\,€:}\quad \text{48\,€ sind 80\,\%} & \\ 1\,\% &= 48\,€ : 80 = 0{,}60\,€ \\ 100\,\% &= 60\,€}
\begin{teile}
\teil Nach einer Senkung um $10\,\%$ kostet ein Pullover $36\,€$. Wie viel kostete er vorher? \leerfeld[€]
\teil Nach einer Erhöhung um $25\,\%$ kostet ein Ticket $15\,€$. Wie viel kostete es vorher? \leerfeld[€]
\teil Ein Gerät kostet $238\,€$ brutto, darin sind $19\,\%$ Mehrwertsteuer enthalten. Wie hoch ist der Nettopreis? \leerfeld[€]
\teil Ein Buch kostet $21{,}40\,€$ brutto, darin sind $7\,\%$ Mehrwertsteuer enthalten. Wie hoch ist der Nettopreis? \leerfeld[€]
\end{teile}
\end{aufgabe}

\begin{aufgabe}{Veränderung – Prozentpunkte. In einer Umfrage stimmen im ersten Jahr $40\,\%$ der Befragten zu, im zweiten Jahr $46\,\%$.}
\begin{teile}
\teil Um wie viele Prozentpunkte ist die Zustimmung gestiegen? \leerfeld
\teil Um wie viel Prozent ist die Zustimmung gestiegen? \leerfeld[\%]
\teil In einer anderen Umfrage sinkt ein Anteil von $25\,\%$ auf $22\,\%$. Gib die Veränderung in Prozentpunkten und in Prozent an. \\ Prozentpunkte: \leerfeld \quad Prozent: \leerfeld[\%]
\end{teile}
\end{aufgabe}

\begin{aufgabe}{Steigung in Prozent. Die Steigung in Prozent gibt an, um wie viele Meter eine Straße auf 100\,m waagerechter Strecke ansteigt.}
\begin{teile}
\teil Eine Straße hat $8\,\%$ Steigung. Das heißt: Auf \leerfeld\ m waagerechter Strecke steigt sie um \leerfeld\ m.
\teil Ein Weg steigt auf $50\,$m waagerechter Strecke um $6\,$m. Wie viel Prozent Steigung sind das? \leerfeld[\%]
\teil Eine Rampe überwindet $21\,$m Höhe auf $250\,$m waagerechter Strecke. Für Radfahrer sind höchstens $6\,\%$ Steigung erlaubt. Prüfe, ob die Rampe zulässig ist. \hfill (P10 2025) \\ Steigung: \leerfeld[\%] \quad zulässig? \janein
\end{teile}
\end{aufgabe}

\begin{aufgabe}{Veränderung – Fehler finden. Der Preis einer Fahrkarte steigt von $60\,€$ auf $75\,€$. Ole rechnet:}
\rechnung{75\,€ - 60\,€ &= 15\,€ \\ 15 : 75 &= 0{,}2 = 20\,\%}
\begin{teile}
\teil Benenne den Fehler und korrigiere die Rechnung. \\ Fehler: \dsleer\dsleer\dsleer \quad richtig: \leerfeld[\%]
\teil Der Preis eines Monatstickets steigt von $45\,€$ auf $54\,€$. Um wie viel Prozent ist er gestiegen? \leerfeld[\%]
\end{teile}
\end{aufgabe}

\begin{aufgabe}{Veränderung – Begründen. Begründe ohne Taschenrechner.}
\begin{teile}
\teil Bei einer Preiserhöhung: Welcher Wert ist $100\,\%$ – der alte oder der neue? \\ \dsleer\dsleer\dsleer\dsleer
\teil Ein Preis wird zuerst um $10\,\%$ erhöht und danach um $10\,\%$ gesenkt. Ist er danach wieder so hoch wie am Anfang? \janein \\ Begründung: \dsleer\dsleer\dsleer\dsleer
\end{teile}
\end{aufgabe}
